\section*{Заключение}
\addcontentsline{toc}{section}{Заключение}

В ходе выполнения курсового проекта разработана платформа \textbf{WhyGame} ---
веб-сервис для публикации и запуска браузерных игр, ориентированный на
инди-разработчиков и студентов. Все поставленные задачи выполнены в полном объёме.

По результатам работы можно сформулировать следующие \textbf{выводы}:

\begin{enumerate}
    \item Стек Flask + PostgreSQL + nginx, упакованный в Docker Compose, оказался
    достаточным для реализации всех функциональных требований платформы при
    минимальных затратах на инфраструктуру. Разделение ответственности между
    сервисами (бизнес-логика, база данных, раздача статики) упрощает
    сопровождение и масштабирование.

    \item Встраивание рейтинга в таблицу комментариев вместо отдельной таблицы
    \texttt{ratings} сократило число таблиц и упростило агрегацию: средний балл
    вычисляется одним SQL-запросом через \texttt{func.avg()} без дополнительных
    JOIN-операций.

    \item Корректная раздача Unity WebGL-сборок потребовала явной настройки nginx:
    каждый тип сжатия (Brotli, Gzip) требует собственного \texttt{location}-блока
    с заголовками \texttt{Content-Encoding}, \texttt{Content-Type},
    \texttt{Cross-Origin-Opener-Policy} и \texttt{Cross-Origin-Embedder-Policy}.
    Без этих заголовков браузер не может ни распаковать файлы, ни активировать
    \texttt{SharedArrayBuffer}.

    \item Двухуровневая система получения роли разработчика (мгновенно) и
    модератора (через заявку с уведомлением) обеспечивает баланс между
    доступностью платформы для новых авторов и контролем над привилегированными
    ролями.
\end{enumerate}

\textbf{Практическая значимость.} Платформа работоспособна и доступна в виде
открытого репозитория (\url{https://github.com/whynotfu/WhyGame}). Любой
разработчик может развернуть собственный экземпляр командой
\texttt{docker compose up --build} при наличии Docker на сервере.

\textbf{Перспективы развития.} В качестве направлений дальнейшей работы можно
выделить: добавление рейтинговой таблицы лидеров по жанрам, поддержку
многофайловых загрузок (Godot, Construct), OAuth-авторизацию через сторонние
сервисы, а также административную панель с агрегированной статистикой по всем
играм платформы.

При оформлении пояснительной записки соблюдались требования~\cite{gost_report}
к структуре и правилам оформления учебных работ.

